\chapter{Безпечний зв'язок та криптографічні протоколи}

\section{Інтеграція та концепція}
Підсистема «ERP/1: Комунікатор» об'єднує в єдиний захищений контур сервіси асинхронного обміну миттєвими повідомленнями (CHAT v2) та поштового обміну документами (MAIL v3). Вона абсорбує технічні рішення:
\begin{itemize}
    \item \textbf{chat.n2o.dev} --- архітектуру WebSocket-брокерів на базі Erlang/OTP;
    \item \textbf{x509.chat} --- автентифікацію за клієнтськими сертифікатами X.509;
    \item \textbf{protocol.zencrypted.uk} --- специфікації криптографічного протоколу Buddha (X.422).
\end{itemize}

\section{Buddha Protocol (Стандарт X.422)}
Протокол Buddha визначає транспортні та криптографічні правила обміну повідомленнями у ненадійних мережах.

\subsection{X.422.1: Транспортний рівень}
Обмін даними здійснюється через асинхронні TCP сокети або шифровані UDP-канали. Для локального виявлення пристроїв у LAN-мережах без DNS-серверів підтримується робота через UDP Multicast на виділений порт 4221. Сервер виконує роль асинхронного брокера, що маршрутизує пакети без збереження їх на сервері після доставки отримувачу.

\subsection{X.422.2: Криптографічний конверт CMS}
Вміст повідомлень та вкладених файлів шифрується на кінцевих пристроях користувачів (End-to-End Encryption) та пакується у CMS-структуру (RFC 5652) у кодуванні ASN.1 DER:
\begin{itemize}
    \item \textbf{SignedData} --- містить підпис ДСТУ 4145 або ECDSA та сертифікат відправника;
    \item \textbf{EnvelopedData} --- містить дані, зашифровані симетричним шифром ДСТУ 7624 (Калина) або AES-GCM-256 на сесійному ключі, який зашифрований на відкритому ключі отримувача за схемою ECDH.
\end{itemize}

\section{Субкомпоненти контуру безпеки}
«Комунікатор» складається з 6 Erlang/OTP додатків:
\begin{enumerate}
    \item \textbf{CA} --- локальний центр випуску сертифікатів за протоколами EST (RFC 7030), CMP (RFC 4210), OCSP статусів та TSP міток часу.
    \item \textbf{VPN} --- клієнт-серверне тунелювання WireGuard на базі сертифікатів відкритих ключів.
    \item \textbf{LDAP} --- ієрархічний довідник користувачів з підтримкою DER-кодування та атрибутики ABAC.
    \item \textbf{IAS} --- служба автентифікації сесій та перевірки статусів відкликання ключів користувачів.
    \item \textbf{CHAT} --- асинхронний WebSocket месенджер v2 з підтримкою Double Ratchet.
    \item \textbf{MAIL} --- сервер асинхронної доставки документів та пошти v3 (SMTP/IMAP over TLS).
\end{enumerate}

\section{Технічна реалізація}

\subsection{Визначення структур даних (Erlang Records)}
Основні сутності підсистеми визначені в Erlang-записах:
\begin{verbatim}
-record(chat_session, {
    session_id,                 % UUID сесії (PK)
    client_id,                  % Серійний номер сертифіката X.509
    device_id,                  % Ідентифікатор пристрою
    ip_address,                 % Поточна IP-адреса
    access_token,               % Токен доступу сесії
    status = active,            % active | expired | revoked
    created_at = 0
}).

-record(ca_certificate, {
    serial_number,              % Серійний номер X.509 (PK)
    subject_dn,                 % Відмітне ім'я суб'єкта
    public_key_der,             % Публічний ключ у DER кодуванні
    not_before = 0,             % Початок дії
    not_after = 0,              % Завершення дії
    status = valid              % valid | revoked | expired
}).
\end{verbatim}

\subsection{Процесний рушій BPE}
Для управління життєвим циклом обробки повідомлень використовується BPMN/FSM опис процесу доставки:
\begin{verbatim}
-module(chat_message).
-include("bpe.hrl").
-compile(export_all).

def() ->
    #process{
        name = 'Buddha Protocol Message Delivery',
        beginEvent = 'ReceiveEnvelope',
        endEvent = 'DestroyMessage',
        tasks = [
            #beginEvent{name='ReceiveEnvelope'},
            #serviceTask{name='RouteByProfile', module=?MODULE},
            #serviceTask{name='StoreTransientQueue', module=?MODULE},
            #serviceTask{name='DeliverToClient', module=?MODULE},
            #serviceTask{name='WaitDeliveryReceipt', module=?MODULE},
            #endEvent{name='DestroyMessage'}
        ],
        flows = [
            #sequenceFlow{name='1', source='ReceiveEnvelope', target='RouteByProfile'},
            #sequenceFlow{name='2', source='RouteByProfile', target='StoreTransientQueue'},
            #sequenceFlow{name='3', source='StoreTransientQueue', target='DeliverToClient'},
            #sequenceFlow{name='4', source='DeliverToClient', target='WaitDeliveryReceipt'},
            #sequenceFlow{name='5', source='WaitDeliveryReceipt', target='DestroyMessage'}
        ],
        roles = [operator, system]
    }.
\end{verbatim}

\subsection{C99 NIF розбирач TLV DER пакетів}
Для швидкого розбору ASN.1 DER об'єктів безпосередньо в Erlang без Cargo/Rust, використовується C99 NIF обгортка:
\begin{verbatim}
#include "erl_nif.h"
#include <string.h>

static ERL_NIF_TERM parse_tlv_nif(ErlNifEnv* env, int argc, const ERL_NIF_TERM argv[]) {
    ErlNifBinary bin;
    if (!enif_inspect_binary(env, argv[0], &bin)) {
        return enif_make_badarg(env);
    }

    if (bin.size < 2) {
        return enif_make_tuple2(env, enif_make_atom(env, "error"), 
                                enif_make_string(env, "too_short", ERL_NIF_LATIN1));
    }

    unsigned char tag = bin.data[0];
    unsigned int length = bin.data[1];
    unsigned int header_len = 2;

    if (length & 0x80) {
        unsigned int num_bytes = length & 0x7F;
        if (bin.size < 2 + num_bytes) {
            return enif_make_tuple2(env, enif_make_atom(env, "error"), 
                                    enif_make_string(env, "invalid_length", ERL_NIF_LATIN1));
        }
        length = 0;
        for (unsigned int i = 0; i < num_bytes; i++) {
            length = (length << 8) | bin.data[2 + i];
        }
        header_len = 2 + num_bytes;
    }

    if (bin.size < header_len + length) {
        return enif_make_tuple2(env, enif_make_atom(env, "error"), 
                                enif_make_string(env, "payload_mismatch", ERL_NIF_LATIN1));
    }

    ERL_NIF_TERM type_term = enif_make_uint(env, tag);
    unsigned char* payload_data;
    ERL_NIF_TERM payload_term;
    payload_data = enif_make_new_binary(env, length, &payload_term);
    memcpy(payload_data, bin.data + header_len, length);

    return enif_make_tuple3(env, enif_make_atom(env, "ok"), type_term, payload_term);
}

static ErlNifFunc nif_funcs[] = {
    {"parse_tlv", 1, parse_tlv_nif}
};

ERL_NIF_INIT(chat_nif, nif_funcs, NULL, NULL, NULL, NULL)
\end{verbatim}
